\section{Budgeted online adaptation: a completed development decision}
\label{sec:adaptive-development}
% Generated numerical macros from the amended adaptive development evidence.
\providecommand{\ADRunCount}{144}
\providecommand{\ADForecastCount}{387{,}900}
\providecommand{\ADAuditChecks}{6{,}713{,}837}
\providecommand{\ADSmallKRRExcess}{22.26}
\providecommand{\ADSmallPninetyNine}{9.89}
\providecommand{\ADLargeKRRExcess}{9.22}
\providecommand{\ADLargePninetyNine}{11.75}
\providecommand{\ADOverruns}{4}
\providecommand{\ADServiceSteps}{21{,}662}
\providecommand{\ADPeakPrediction}{175{,}555.08}
\providecommand{\ADPeakTruth}{-0.9815}
\providecommand{\ADQueryFeatureMagnitude}{3.20}
\providecommand{\ADReplayFeatureMagnitude}{392}
\providecommand{\ADRelativeDenominator}{0.354}


\paragraph{Question and scope.}
A subsequent development study tested whether updates to local KAHM geometry under fixed reference and
replay budgets provide enough predictive benefit to justify a larger evaluation.
The complete design comprises twelve streams, six methods and two capacity
tiers: \ADRunCount{} runs and \ADForecastCount{} scored forecasts after the
numerical amendment described below. Four stationary nonlinear AR streams
are paired by innovations with four recurring A/B streams; the latter
alternate regimes every 256 observations, and both regimes occur in the
initial fitting prefix. Regime labels are never supplied to the forecasters.
Three Dominick weekly series and one Sunspot daily series are deliberately
reused from previously inspected inputs. This supplies development evidence
from one generator and two selected natural collections, with no claim of
new-source confirmation. Appendix~\ref{app:adaptive-development} records
the chronology, decision thresholds and source identities.

\paragraph{A simpler adaptive prototype.}
All geometric methods retain a sliding-window raw ridge baseline and
residual heads. Fixed KAHM updates its heads with unchanged initial geometry;
adaptive KAHM, affine and RBF-local share the bounded reference-refresh
policy. Each adaptive refresh coherently re-encodes replay histories and
rebuilds residual heads. The sixth method is sliding-window Gaussian kernel
ridge with a free intercept and recursive add/delete inverse updates. Its
forecast also uses the calibrated blend with raw ridge: correction weight
one gives the kernel prediction and zero gives raw ridge. Thus the reported
KRR arm is a selected forecasting procedure, not an unblended kernel forecast
on every horizon, and it is not the published KRLS-T algorithm.
The small tier allows 64 shared
geometric references and 128 replay rows per horizon, and the large tier
allows 128 and 256. Kernel supports are capped separately per horizon;
these quotas are not equal byte or computational allocations. Bandwidths
and gate temperatures are fixed, so this prototype is not a streaming
reproduction of the forecast-trained model in the preceding study.

\paragraph{Chronology and numerical amendment.}
The first 35\% of each stream initializes the model, 35--50\% supplies
chronological calibration, and the latter half supplies evaluation.
Only matured labels enter an update. Forecasts are recorded before their
targets arrive and remain immutable. Each horizon selects a correction
weight from $\{0,0.25,0.5,1\}$ on matured calibration targets, then retains
that weight during evaluation. The original complete attempt had four
correction-head failures and 8,659 framework fallback forecasts. Amendment
A1 replaces accumulated correction normal equations with a centered,
column-equilibrated augmented ridge solve using pivoted QR, preserving
the objective, penalties, data, scalers, reference policy and decision
rules. All geometric arms receive the repair, and the full roster is
repeated. The amended runs have no unexpected model exceptions, framework
fallback forecasts or observed cap violations. Both attempts are preserved;
their source freezes were local checkpoints, not public preregistration.
The repeat follows observed outcomes and is not independent confirmation.

% Generated from saved amended development evidence; no model fitting.
\begin{table}[!htbp]
\centering\small\setlength{\tabcolsep}{5pt}
\caption{Budgeted adaptive development study: scaled MSE, lower is better. Origins are averaged within horizon, horizons within series, and series within collection. The natural mean gives Dominick and Sunspot equal weight. The two capacity tiers are quotas, not equal-byte allocations. All non-ridge arms, including KRR, use calibration-selected mixing with the common ridge baseline. No losses are pooled with earlier studies.}
\label{tab:adaptive-development-loss}
\begin{tabular}{@{}lrrrrr@{}}
\toprule
Method & Stationary & Recurring & Dominick & Sunspot & Natural mean \\
\midrule
\multicolumn{6}{l}{\textit{Small: 64 shared references / 128 replay rows per horizon}} \\
Raw ridge & 2.2402 & 3.4583 & 1.6144 & 2.3697 & 1.9920 \\
Fixed KAHM & 2.2407 & 3.4583 & 1.6134 & 2.3697 & 1.9915 \\
Adaptive KAHM & 2.2399 & 3.4583 & 1.6135 & 2.3697 & 1.9916 \\
Adaptive affine & 2.2407 & 3.4583 & 1.6137 & 3.4004 & 2.5071 \\
Adaptive RBF-local & 2.2428 & 3.4583 & 1.6160 & 2.3697 & 1.9928 \\
Sliding-window KRR & 2.2491 & 3.3578 & 0.8685 & 2.3897 & 1.6291 \\
\midrule
\multicolumn{6}{l}{\textit{Large: 128 shared references / 256 replay rows per horizon}} \\
Raw ridge & 2.0275 & 3.3623 & 0.8471 & 2.1045 & 1.4758 \\
Fixed KAHM & 2.0375 & 3.4031 & 0.8464 & 2.1045 & 1.4754 \\
Adaptive KAHM & 2.0280 & $1.376\times 10^{8}$ & 0.8444 & 2.1045 & 1.4744 \\
Adaptive affine & 2.0364 & $1.061\times 10^{6}$ & 0.8446 & 2.2085 & 1.5266 \\
Adaptive RBF-local & 2.0290 & 3.7036 & 0.8460 & 2.1045 & 1.4752 \\
Sliding-window KRR & 2.0551 & 3.2079 & 0.5954 & 2.1045 & 1.3499 \\
\bottomrule
\end{tabular}
\end{table}


\paragraph{Predictive outcome.}
Table~\ref{tab:adaptive-development-loss} reports MSE divided by an
initial-prefix seasonal-difference scale. The scale and online schedule
differ from earlier sections, precluding direct comparisons of their
numerical scores. Both complete tiers fail the recurring and natural-data
accuracy criteria and receive the decision \textsc{Defer}. Adaptive KAHM's
natural mean loss exceeds kernel ridge by \ADSmallKRRExcess\% at the small
tier and \ADLargeKRRExcess\% at the large tier. At each tier, only two of
twelve natural horizon/series combinations activate a KAHM correction;
every Sunspot horizon selects zero. Small-tier recurring corrections are
all zero, while nine of twelve large-tier recurring corrections are active.
Baseline recovery at zero correction does not demonstrate added geometric
value. The descriptive advancement criteria are development triage rules,
not a hypothesis test or a universal definition of effectiveness.

\paragraph{Extrapolation persists after the numerical repair.}
For recurring seed 20261101, origin 1554 and horizon 12, large-tier adaptive
KAHM predicts $\ADPeakPrediction$ when the target is $\ADPeakTruth$.
At the preserved diagnostic state, the largest absolute standardized query
feature is about $\ADQueryFeatureMagnitude\times10^6$, compared with
$\ADReplayFeatureMagnitude$ in the current supervised replay. A fresh
centered SVD solve reproduces the large correction. The minimum relative
signed denominator, defined by the ratio in
Equation~\eqref{eq:mechanism-relative-denominator}, is
$\ADRelativeDenominator$, safely above its $10^{-10}$
guard. Adaptive affine also exhibits severe extrapolation under a fresh
solve. These post-run observations identify a representation--decoder
stability problem in the examined states. They are consistent with a
query falling outside the observed feature range of matured-label replay, but do not
isolate the causal contributions of scaling, label delay and geometry
refreshes. No forecasts were removed, clipped or retuned to obtain the
reported decision.

\paragraph{Cost and research consequence.}
Adaptive KAHM passes the monitored 20\,ms p99 screen on every stream:
the worst per-stream values are \ADSmallPninetyNine{} and
\ADLargePninetyNine{}\,ms. There are nevertheless \ADOverruns{} service
intervals above 20\,ms among \ADServiceSteps{} evaluation observation
steps; the screen is not an enforced per-step deadline. Service includes
updates, geometry refreshes, feature/head rebuilds and prediction, but
excludes evidence serialization and scoring. The eager prototype maintains
geometry even when its correction is disabled. Resource details appear
in Table~\ref{tab:adaptive-development-resources}.
This result establishes measured implementation feasibility under the
stated screen while withholding support for an adaptive forecasting
advantage. The fixed-geometry reconstruction-departure experiment in
Section~\ref{sec:reconstruction-mechanism} subsequently tests a distinct
readout and an explicit correction bound. Its result is reported separately,
rather than treated as a repair or confirmation of the adaptive prototype.
The development outcome does not rule out every other adaptive or
forecast-trained KAHM formulation.
